%% ma: L10-B15-0108
%% lop: 10
%% chuong: 6
%% bai: 15
%% dang: 10.15.4
%% loai: DS
%% muc: [NB, NB, NB, TH]
%% dap_an: "SĐĐS"
%% nguon: "Ảnh giáo viên gửi 25/11/2026 (lô 2): Bài tập trắc nghiệm – Toán 10 Bài 15 Hàm số (sự biến thiên, đồ thị), Phần II Câu 4"
%% kien_thuc: "Lớp 10 Bài 15 (đồ thị, tập giá trị, đồng biến nghịch biến)"
%% trang_thai: chinh-thuc
%% ghi_chu: "Đồ thị vẽ lại theo y = -x^2/2 (khớp các điểm (-1;-1/2), (-2;-2) của hình gốc)."
\begin{ex}
Cho hàm số $f(x)$ có đồ thị như hình vẽ. Xét tính đúng sai của các khẳng định sau:
\begin{center}\begin{tikzpicture}[scale=.9]
\draw[truc] (-3,0)--(3,0) node[below]{$x$};
\draw[truc] (0,-3.4)--(0,.8) node[left]{$y$};
\draw[phu] (-2,0)--(-2,-2)--(0,-2) (-1,0)--(-1,-.5)--(0,-.5);
\node[above] at (-2,0) {\footnotesize$-2$}; \node[above] at (-1,0) {\footnotesize$-1$};
\node[above right] at (0,0) {\footnotesize$O$};
\node[right] at (0,-.5) {\footnotesize$-\frac12$}; \node[right] at (0,-2) {\footnotesize$-2$};
\draw[thick,domain=-2.5:2.5,smooth] plot (\x,{-(\x)^2/2});
\end{tikzpicture}\end{center}
\choiceTF
{Điểm $A\left(-1;\dfrac12\right)$ thuộc đồ thị của hàm số.}
{\True Hàm số $f(x)$ nghịch biến trên khoảng $(0;+\infty)$.}
{\True Hàm số $f(x)$ đồng biến trên khoảng $(-2;0)$.}
{Tập giá trị của hàm số là $\mathbb R$.}
\loigiai{a) $f(-1)=-\tfrac12\ne\tfrac12$: Sai. b) Đồ thị đi xuống trên $(0;+\infty)$: Đúng. c) Đồ thị đi lên trên $(-\infty;0)$ nên đồng biến trên $(-2;0)$: Đúng. d) $f(x)\le0$ nên tập giá trị là $(-\infty;0]$: Sai.}
\end{ex}
